\documentclass{report}
\usepackage{pdfpages}
\usepackage{graphicx} 
\usepackage{dirtytalk}
\usepackage{lmodern}
\usepackage{booktabs}
\usepackage{listings}
\usepackage{float}
\usepackage{siunitx}
\usepackage[hidelinks]{hyperref}
\usepackage{amsmath}
\usepackage{rotating}
\usepackage{arydshln}
\usepackage[ngerman]{babel}
\usepackage[a4paper, total={6in, 8in}]{geometry}
\begin{document}
\begin{titlepage}
    \begin{center}
        \vspace*{1cm}
            
        \Huge
        \textbf{Auswertung und Protokoll - Korrektur vom 08.03.2026}
            
        \vspace{0.5cm}
        \LARGE
        Zum Versuch: MOS - Mechanischer Oszillator
        
        \vspace{1.8cm}

        \textbf{Jonas Müther \& Alejandro Schultheiss}
        \vfill
            
        P1 Praktikum\\
       
            
        \vspace{0.8cm}
  

        \vspace{1.8cm}
        \Large
        LMU München \\
        Physik B.Sc.\\
        Deutschland\\
        2026
            
    \end{center}
\end{titlepage}
\tableofcontents


\newpage
\section{Korrektur zu Kommentar 25.1}
Somit sind die Endergebnisse für die Frequenzen:
\begin{equation}
    \begin{split}
        f_{m_1} = (1.3928 \pm 0.0005) \si{\hertz} \\
        f_{m_2} = (1.3021 \pm 0.0006) \si{\hertz} \\
        f_{m_3} = (1.22549 \pm 0.00025) \si{\hertz} \\
        f_{m_4} = (1.1641 \pm 0.0004) \si{\hertz} \\
        f_{m_5} = (1.1161 \pm 0.0006) \si{\hertz} \\
    \end{split}
\end{equation}




\section{Korrektur zu Kommentar 29.1}
\[
\Delta f_d
=
\frac{1}{\bar T^2}\,\Delta \bar T
=
\frac{0{,}0123}{(0{,}8744)^2}\,\mathrm{Hz}
\approx 0{,}016\,\mathrm{Hz}.
\]

\noindent
Da die führende signifikante Ziffer eine 1 ist, wird die Unsicherheit mit zwei signifikanten Stellen angegeben.
Damit folgt

\[
\boxed{f_d = (1{,}144 \pm 0{,}016)\,\mathrm{Hz}}.
\]




\section{Korrektur zu Kommentar 30.1}
Theoretisch ist die Frequenz der freien gedämpften Schwingung kleiner als die der ungedämpften Schwingung, da
\[
\omega_\mathrm d=\sqrt{\omega_0^2-\beta^2}
\quad\Rightarrow\quad
f_\mathrm d=\frac{\omega_\mathrm d}{2\pi}<\frac{\omega_0}{2\pi}=f_0.
\]
Die Frequenzen müssen daher nicht identisch sein; vielmehr ist $f_\mathrm d<f_0$ zu erwarten.
Im Experiment ergibt sich jedoch $f_\mathrm d>f_0$. Dies deutet auf einen dominierenden systematischen bzw. statistischen Einfluss in der Periodenbestimmung (z.\,B.\ geringe Anzahl an Perioden, abklingende Amplitude/Triggerung) hin.
Unter Berücksichtigung der Unsicherheiten beträgt die Abweichung etwa $1{,}7\sigma$ und liegt damit zwar noch im Bereich einer $2\sigma$-Verträglichkeit, die Vorzeichenrichtung widerspricht jedoch der theoretischen Erwartung.

\section{Korrektur zu Kommentar 31.1}

Nach Gleichung (32) gilt
\[
\Lambda=\beta T_d,
\qquad\Rightarrow\qquad
\beta=\frac{\Lambda}{T_d}.
\]
\noindent
Mit
\[
\Lambda=0{,}862,
\qquad
T_d=\bar T=(0{,}8744\pm 0{,}0123)\,\mathrm{s}
\]
ergibt sich für den Zentralwert
\[
\beta=\frac{0{,}862}{0{,}8744\,\mathrm{s}}=0{,}9858\,\mathrm{s^{-1}}.
\]

Die Unsicherheit wird mittels Gaußscher Fehlerfortpflanzung bestimmt:
\[
u(\beta)=\sqrt{\left(\frac{1}{T_d}u(\Lambda)\right)^2+
\left(\frac{\Lambda}{T_d^2}u(T_d)\right)^2}.
\]
Da keine Unsicherheit für $\Lambda$ vorliegt, wird $u(\Lambda)\approx 0$ gesetzt, sodass
\[
u(\beta)\approx \frac{\Lambda}{T_d^2}\,u(T_d)
=\frac{0{,}862}{(0{,}8744)^2}\cdot 0{,}0123\,\mathrm{s^{-1}}
\approx 0{,}0139\,\mathrm{s^{-1}}.
\]

Gerundet folgt damit
\[
\boxed{\beta=(0{,}986\pm 0{,}014)\,\mathrm{s^{-1}}.}
\]
\section{Korrektur zu Kommentar 32.1}

Nach Gleichung (34) gilt
\[
T_d = T_0\sqrt{1+\left(\frac{\Lambda}{2\pi}\right)^2},
\qquad
T_0=\frac{1}{f_0}.
\]
\noindent
Mit
\[
f_0=(1{,}1161\pm 0{,}0006)\,\mathrm{Hz}
\]
\noindent
ergibt sich
\[
T_0=\frac{1}{f_0}=0{,}89598\,\mathrm{s}.
\]
\noindent
Die Unsicherheit von $T_0$ folgt aus Gaußscher Fehlerfortpflanzung:
\[
u(T_0)=\left|\frac{\partial}{\partial f_0}\frac{1}{f_0}\right|u(f_0)
=\frac{1}{f_0^2}\,u(f_0)
=\frac{0{,}0006}{(1{,}1161)^2}\,\mathrm{s}
=0{,}000482\,\mathrm{s}.
\]
\noindent
Setzt man
\[
F=\sqrt{1+\left(\frac{\Lambda}{2\pi}\right)^2},
\]
so ist $T_d=T_0F$. Mit $\Lambda=0{,}862$ folgt
\[
F=\sqrt{1+\left(\frac{0{,}862}{2\pi}\right)^2}=1{,}00937,
\]
\[
T_d = 0{,}89598\cdot 1{,}00937 = 0{,}90437\,\mathrm{s}.
\]
\noindent
Da keine Unsicherheit für $\Lambda$ vorliegt, wird $u(\Lambda)$ vernachlässigt. Damit ergibt sich
\[
u(T_d)\approx \left|\frac{\partial T_d}{\partial T_0}\right|u(T_0)=F\,u(T_0)
=1{,}00937\cdot 0{,}000482\,\mathrm{s}
=0{,}000486\,\mathrm{s}.
\]
\noindent
Somit erhält man
\[
\boxed{T_d^{(\mathrm{theo})}=(0{,}9044\pm 0{,}0005)\,\mathrm{s}.}
\]




\section{Korrektur zu Kommentar 32.2}

Es wurde bestimmt:
\[
T_d^{(\mathrm{theo})}=(0{,}9044\pm 0{,}0005)\,\mathrm{s},
\qquad
\bar T_d^{(\mathrm{exp})}=(0{,}8744\pm 0{,}0123)\,\mathrm{s}.
\]

Die absolute Abweichung beträgt:
\[
\Delta T = T_d^{(\mathrm{theo})}-\bar T_d^{(\mathrm{exp})}
=0{,}9044-0{,}8744
=0{,}0300\,\mathrm{s}.
\]

Die relative Abweichung ergibt sich zu:
\[
\frac{\Delta T}{\bar T_d^{(\mathrm{exp})}}
=\frac{0{,}0300}{0{,}8744}
\approx 0{,}0343 \approx 3{,}4\%.
\]

Zur Beurteilung der Übereinstimmung wird die kombinierte Unsicherheit der Differenz verwendet:
\[
u_{\Delta T}=\sqrt{u(T_d^{(\mathrm{theo})})^2+u(\bar T_d^{(\mathrm{exp})})^2}
=\sqrt{(0{,}0005)^2+(0{,}0123)^2}
\approx 0{,}0123\,\mathrm{s}.
\]

Damit folgt die normierte Abweichung:
\[
z=\frac{|\Delta T|}{u_{\Delta T}}
=\frac{0{,}0300}{0{,}0123}
\approx 2{,}4.
\]

Die Abweichung liegt damit bei etwa $2{,}4\sigma$, also geringfügig oberhalb eines $2\sigma$-Bereichs.
Somit zeigen Theorie und Experiment eine insgesamt gute Größenordnung, jedoch keine vollständige Übereinstimmung innerhalb eines strengen $2\sigma$-Kriteriums.
Da keine Unsicherheit für $\Lambda$ vorliegt, ist die theoretische Unsicherheit hierbei vermutlich unterschätzt.
\section{Korrektur zu Kommentar 34.1}
\subsubsection*{Einordnung: Soll $f_\mathrm{res}$ mit $f_0$ bzw. $f_\mathrm{d}$ übereinstimmen?}

Die Resonanzfrequenz $f_\mathrm{res}$ (Maximum der Amplitude der erzwungenen Schwingung) muss theoretisch nicht mit der Eigenfrequenz der freien ungedämpften Schwingung $f_0$ übereinstimmen.
Für den gedämpft erzwungenen Oszillator gilt nach Gleichung (40)

\[
\omega_\mathrm{res}=\sqrt{\omega_0^2-2\beta^2}.
\]

Da für $\beta>0$ der Term $2\beta^2$ positiv ist, folgt unmittelbar

\[
\omega_\mathrm{res}<\omega_0
\qquad\Rightarrow\qquad
f_\mathrm{res}<f_0.
\]

Zudem gilt für die freie gedämpfte Schwingung $\omega_\mathrm{d}=\sqrt{\omega_0^2-\beta^2}$, womit auch
\[
\omega_\mathrm{res}<\omega_\mathrm{d}
\qquad\Rightarrow\qquad
f_\mathrm{res}<f_\mathrm{d}
\]
zu erwarten ist.

Im Experiment ergibt sich $f_\mathrm{res}=1{,}113\,\mathrm{Hz}$ und damit tatsächlich $f_\mathrm{res}<f_0$.
Die Resonanzfrequenz liegt also in der erwarteten Richtung unterhalb der ungedämpften Eigenfrequenz.
Eine vollständige numerische Übereinstimmung ist nicht zu erwarten, da Dämpfung die Resonanzfrequenz verschiebt.


\section{Korrektur zu Kommentar: 35.1}
\subsubsection*{Berechnung der Resonanzkreisfrequenz $\omega_\mathrm{res}$ (mit Unsicherheit)}

Nach Gleichung (40) gilt
\[
\omega_\mathrm{res}=\sqrt{\omega_0^2-2\beta^2}.
\]

Zunächst wird aus $f_0$ die Kreisfrequenz $\omega_0$ bestimmt:
\[
\omega_0=2\pi f_0.
\]
Mit
\[
f_0=(1{,}1161\pm 0{,}0006)\,\mathrm{Hz}
\]
folgt
\[
\omega_0=2\pi\cdot 1{,}1161=7{,}0127\,\mathrm{rad/s}.
\]
Die Unsicherheit von $\omega_0$ ergibt sich zu
\[
u(\omega_0)=2\pi\,u(f_0)=2\pi\cdot 0{,}0006=0{,}00377\,\mathrm{rad/s}.
\]

Für die Dämpfungskonstante gilt
\[
\beta=(0{,}986\pm 0{,}014)\,\mathrm{s^{-1}}.
\]

Damit ergibt sich
\[
\omega_\mathrm{res}=\sqrt{(7{,}0127)^2-2\,(0{,}986)^2}=6{,}8726\,\mathrm{rad/s}.
\]

Die Unsicherheit wird durch Gaußsche Fehlerfortpflanzung bestimmt. Mit
\[
\frac{\partial \omega_\mathrm{res}}{\partial \omega_0}=\frac{\omega_0}{\omega_\mathrm{res}},
\qquad
\frac{\partial \omega_\mathrm{res}}{\partial \beta}=-\frac{2\beta}{\omega_\mathrm{res}}
\]
folgt
\[
u(\omega_\mathrm{res})=
\sqrt{\left(\frac{\omega_0}{\omega_\mathrm{res}}u(\omega_0)\right)^2+
\left(\frac{2\beta}{\omega_\mathrm{res}}u(\beta)\right)^2 }.
\]
Einsetzen liefert
\[
u(\omega_\mathrm{res})\approx
\sqrt{\left(\frac{7{,}0127}{6{,}8726}\cdot 0{,}00377\right)^2+
\left(\frac{2\cdot 0{,}986}{6{,}8726}\cdot 0{,}014\right)^2}
\approx 0{,}0056\,\mathrm{rad/s}.
\]

Somit ergibt sich
\[
\boxed{\omega_\mathrm{res}=(6{,}873\pm 0{,}006)\,\mathrm{rad/s}}.
\]

Optional kann die zugehörige Resonanzfrequenz in Hertz angegeben werden:
\[
f_{\mathrm{res,theo}}=\frac{\omega_\mathrm{res}}{2\pi}
=\frac{6{,}8726}{2\pi}
=1{,}094\,\mathrm{Hz},
\qquad
u(f_{\mathrm{res,theo}})=\frac{u(\omega_\mathrm{res})}{2\pi}\approx 0{,}0009\,\mathrm{Hz}.
\]
Damit folgt
\[
\boxed{f_{\mathrm{res,theo}}=(1{,}094\pm 0{,}001)\,\mathrm{Hz}}.
\]

\section{Korrektur zu Kommentar 35.2}

\subsubsection*{Diskussion der Abweichungen der Dämpfungskonstante}

Die beiden aus den Resonanzkurven bestimmten Werte $\beta_{\text{Amp}}$ und $\beta_{\text{Phase}}$ stimmen mit einer Abweichung von nur etwa $3{,}6\%$ gut miteinander überein. Dies ist plausibel, da beide Größen auf derselben Messreihe der erzwungenen Schwingung beruhen und somit ähnliche systematische Einflüsse teilen. 

Im Vergleich dazu liegt der aus der gedämpften freien Schwingung bestimmte Wert $\beta_{\text{frei}}$ um etwa $15$--$17\%$ höher. Mögliche Ursachen für diese Abweichung sind:

\begin{itemize}
    \item \textbf{Unterschiedliche Auswertungsmethoden und Gewichtung der Messwerte:}
    Bei der freien Schwingung wird $\beta$ aus dem exponentiellen Abklingen der Amplitude bestimmt (über $\Lambda$). Einzelne frühe Messpunkte mit großer Amplitude beeinflussen das Ergebnis dabei stark. Die Resonanzkurven-Fits (Amplitude/Phase) nutzen hingegen viele Frequenzpunkte und liefern dadurch oft stabilere Fitparameter.
    
    \item \textbf{Nichtideale Dämpfung (nicht rein viskos):}
    Die Theorie setzt eine lineare Dämpfung proportional zur Geschwindigkeit voraus. In der Praxis können zusätzlich Reibungseffekte (z.\,B.\ Lagerreibung bzw.\ Coulomb-Reibung) oder amplitudenabhängige Verluste auftreten. Solche Effekte beeinflussen besonders den freien Abklingvorgang und können zu einem scheinbar größeren $\beta_{\text{frei}}$ führen.
    
    \item \textbf{Amplitudenabhängige oder frequenzabhängige Verluste:}
    Bei großer Auslenkung (zu Beginn der freien Schwingung) können zusätzliche Dissipationsmechanismen stärker wirken als bei kleinen Amplituden bzw.\ im stationären erzwungenen Betrieb. Dadurch kann $\beta_{\text{frei}}$ größer erscheinen als die im stationären Zustand bestimmten $\beta$-Werte.
    
    \item \textbf{Mess- und Ableseefekte bei der freien Schwingung:}
    Die Bestimmung der Maxima (und damit $\Lambda$) ist empfindlich gegenüber Ablese- bzw.\ Triggerungenauigkeiten, insbesondere wenn die Amplitude schnell abklingt und das Signal nahe am Rauschen liegt. Dies kann die Steigung der linearen Ausgleichsgeraden und damit $\Lambda$ und $\beta_{\text{frei}}$ systematisch verschieben.
    
    \item \textbf{Leicht unterschiedliche Versuchsbedingungen:}
    Kleine Änderungen in der Aufhängung, Justage oder Dämpfereinstellung zwischen Teilversuch 3 (frei) und Teilversuch 4 (erzwungen) können die effektive Dämpfung beeinflussen und somit zu unterschiedlichen $\beta$-Werten führen.
\end{itemize}

Insgesamt ist die gute Übereinstimmung von $\beta_{\text{Amp}}$ und $\beta_{\text{Phase}}$ ein Hinweis auf eine konsistente Auswertung der erzwungenen Schwingung. Die größere Abweichung zu $\beta_{\text{frei}}$ lässt sich plausibel durch nichtideale Dämpfungsmechanismen sowie durch die höhere Empfindlichkeit der freien Abklingauswertung gegenüber Mess- und Ableseeinflüssen erklären.


\end{document}
